\section{The open problem: selection}\label{sec:selection}

\paragraph{The question.}
Section~\ref{sec:boundary} located what the structure leaves free. Here we list the free ``which'' choices, state the
shape that a principle selecting among them would have to take, and stop. We do not propose such a principle. Every
statement in this section that is not a cited result of an earlier section is a question.

\paragraph{Every ``which'' is free.} Each of the following is supplied, not derived, in every frame on record.
\begin{itemize}
\item The modulus $\tau$: with the Higgs along $\rho$, the record's invariant quartic selects $\tau=\omega$ or runs to the
  cusp. Neither, as a point, carries a hierarchy (B1652); toward the cusp of moduli the first generation's lift falls
  exponentially (B1667). The record's check covers that direction only.
\item The direction of the Higgs in the triplet's space: it is the second-to-third mass matrix (a bijection onto the
  symmetric $2\times2$ matrices), and no constant direction is covariant under the moves (Section~\ref{sec:couplings}).
\item The line, that is the flat hypercharge line of a sector, among the eighteen admissible lines at $\tau=\omega$
  (B1660).
\item The vacuum, and the frame (Verdict~\ref{bd:tribunal}).
\item The side of the end at which the motion settles (B1663): the pull-back's strength sets the probability of the side and the speed, and
  the choice of side is the positivity of the earlier foundation, now read as a choice.
\item The choice of canonical horocycle at the puncture, which fixes the end radius $\mathcal E$ at order one; the choice is
  discrete, and the record does not make it (Census law~\ref{cl:cusp}, B1667).
\end{itemize}
These choices are not counted in the record's price, four axioms and nine irreducible identifications for the whole
chain. No residual symmetry of the weave's group that reduces the count is compatible with the CKM data at $3\sigma$
(Theorem~\ref{bd:reducesnone}).

\paragraph{The shape of a selection principle.}
Two proved statements constrain it. First, an invariant selector fixes an orbit, not a point of it (B990), and the
symmetric points of the modulus carry rank-one couplings or a hierarchy set by data the symmetry does not fix (B1653,
B1657). Symmetry and orbit invariants alone therefore cannot select a point that is not fixed; an invariant dynamics with a
symmetry-breaking vacuum is not excluded. A selection principle must act where the symmetry is broken or at
its fixed points. Second, the clock and the tick supply a grading and a scalar and nothing else (Theorem~\ref{bd:clock}),
and the pull-back's strength supplies a side and a speed but no size (Census law~\ref{bd:strength}). A principle that
takes its input from either is not enough. What would be enough is stated by the demands of the tribunal: a dynamics on
record and run that selects a vacuum, with the completion and the normalisation named among its inputs
(Verdict~\ref{bd:tribunal}).

\paragraph{The measurer inside the weave.}
A measured value is a comparison between a part of the world and a part that keeps a record of it. Two statements bound
what counting the measurer inside the weave can supply. Neither selects anything.

\begin{lemma}[composition]\label{lem:composition}
Let $r(\gamma,\gamma')>0$ be a read-out on pairs of threads that is to be a ratio of measured magnitudes, so that
$r(\gamma,\gamma)=1$ and $r(\gamma,\gamma'')=r(\gamma,\gamma')\,r(\gamma',\gamma'')$. Then
$r(\gamma,\gamma')=f(\gamma)/f(\gamma')$, with $f(\gamma)=r(\gamma,\gamma_0)$ for any fixed thread $\gamma_0$.
\end{lemma}
\begin{proof}
$r(\gamma,\gamma')=r(\gamma,\gamma_0)\,r(\gamma_0,\gamma')$, and $r(\gamma_0,\gamma')=1/r(\gamma',\gamma_0)$.
\end{proof}
\noindent So, for magnitudes read as composing ratios, a measurer among the threads adds one degree of freedom, the unit.
What remains is a function $f$ on the threads, the form of the per-field read-outs (B1668). The lemma says nothing about
what fixes $f$: an interaction among the threads could constrain it, and read-outs that are not scalar ratios are
outside its scope. A pairwise invariant that does not
compose cannot be a ratio of masses: an example is the linking number of two threads as modular
knots~\cite{Ghys2007,DukeImamogluToth2017}. It can stand only for relational data, such as the mixing between two sectors.

\begin{proposition}[no thread keeps a record, B1675]\label{prop:records}
Take the common point, the quaternion representation $\rho_Q$ ($a\mapsto i$, $b\mapsto j$), with coefficients in the
blocks $\mathrm{Sym}^{2k}\rho_Q$ of $\mathfrak{sl}(n)=\bigoplus_{k<n}\mathrm{Sym}^{2k}$. The private states are the kernel
of the restriction $H^1(F_2;\mathrm{Sym}^{2k}\rho_Q)\to H^1$ of the puncture, $z\mapsto z([a,b])$: the classes that the
puncture, a commutator, cannot detect. They are the flat twists of the block's $Q_8$-invariant part, and at rank three
($k=1,2$) they span four dimensions. On them:
\begin{itemize}
\item no thread keeps a non-zero vector, on all $758$ states to word length $12$;
\item the joint action of $L$ and $R$ keeps none either;
\item a single shear can keep one, and so can the involution $P$: at rank three $L$ and $R$ each keep one class, and
  $P$ keeps two (three at rank four).
\end{itemize}
\end{proposition}
\noindent\emph{Proof and scope.} A thread acts on the private states as its own hyperbolic matrix (eigenvalues
$\lambda^{\pm1}$, $\lambda>1$) tensored with a map of finite order, so it has no fixed vector; a shear is unipotent and
has one. The tables at the common point are exact. The proposition concerns fixed vectors. A relation between two
changing parts can persist without either being fixed (if $x\mapsto Ax$ and $\ell\mapsto\ell A^{-1}$, then $\ell x$ is
kept), so it does not exclude a relational record; whether one exists between the private states and a measurer's
states is a question about equivariant maps between them, not computed here. The census over the threads was decided at $60$ digits on $747$ of
the $758$ states, and on the other $11$ at $120$ digits after the run (B1675).

\reading{The motion's facts, as fenced outlook. The layers' normalised coherence climbs toward one along the cusp family, and not at the root $LR$ (Scoreboard~\ref{bd:layers}). Whether that
selects anything depends on the physical amplitude, which this paper does not construct; we leave the question as stated.
Two computations bound the obvious attempts: in the surface frame's twisted units on the least-volume closing threads, and
in the per-field thread frame at every depth computed, no assignment of the charges, the Standard Model's or a null's,
passes the six-ratio mass test, and the hierarchies come out too strong (B1668).}

\reading{The weave and the wave are the two sides of one equation. The Selberg trace formula equates a sum over the
closed geodesics (the threads, by length) with a sum over the spectrum of the modular surface (the waves, by
frequency), plus terms for the identity, the elliptic points and the cusp; the geometric side runs over the conjugacy
classes that the grammar, with the pull-back, makes reachable states. The shape a selection principle would have to take
is one that reads one side through the other. The layers' constant mode (Scoreboard~\ref{bd:layers}) is this equation's
constant term seen on the program's own quantity. That physics lives on this is not shown.}
